%% =============================================================================
%% Glossary Definitions
%% Parkinson's Disease Voice Classification — MSc Thesis
%%
%% Organisation:
%%   Part A — Abbreviations & Acronyms  (\newabbreviation)  — 27 entries
%%   Part B — Technical Terms           (\newglossaryentry)  —  6 entries
%%                                                     Total: 33 entries
%%
%% Abbreviation style: short-only  →  \gls{key} always prints the short form.
%% All entries are forced into the printed glossary via \glsaddall in main.tex,
%% so no entry needs a \gls{} call to appear on the glossary pages.
%% =============================================================================

%% -----------------------------------------------------------------------------
%% PART A — ABBREVIATIONS (type = abbreviations)
%% Format:  \newabbreviation{key}{SHORT}{Ελληνική πλήρης μορφή (English long form)}
%% -----------------------------------------------------------------------------

%% — Core Clinical / Disease ——————————————————————————————————————————————————

\newabbreviation{pd}{PD}{Νόσος Πάρκινσον (Parkinson's Disease)}

\newabbreviation{hc}{HC}{Υγιείς Μάρτυρες Ελέγχου (Healthy Controls)}

\newabbreviation{hkd}{HKD}{Υποκινητική Δυσαρθρία (Hypokinetic Dysarthria)}

\newabbreviation{updrs}{UPDRS}{Ενιαία Κλίμακα Αξιολόγησης Νόσου Πάρκινσον
    (Unified Parkinson's Disease Rating Scale)}

%% — Machine Learning — Models ————————————————————————————————————————————————

\newabbreviation{ml}{ML}{Μηχανική Μάθηση (Machine Learning)}

\newabbreviation{svm}{SVM}{Μηχανή Διανυσμάτων Υποστήριξης (Support Vector Machine)}

\newabbreviation{rbf}{RBF}{Συνάρτηση Ακτινικής Βάσης (Radial Basis Function)}

\newabbreviation{rf}{RF}{Τυχαίο Δάσος (Random Forest)}

\newabbreviation{gb}{GB}{Ενίσχυση Κλίσης (Gradient Boosting)}

\newabbreviation{xgb}{XGB}{Extreme Gradient Boosting (XGBoost)}

\newabbreviation{logr}{LR}{Λογιστική Παλινδρόμηση (Logistic Regression)}

%% — Evaluation Metrics & Validation ————————————————————————————————————————

\newabbreviation{roc}{ROC}{Χαρακτηριστική Λειτουργία Δέκτη
    (Receiver Operating Characteristic)}

\newabbreviation{auc}{AUC}{Εμβαδόν κάτω από την Καμπύλη (Area Under the Curve)}

\newabbreviation{rocauc}{ROC-AUC}{Εμβαδόν κάτω από την Καμπύλη ROC
    (Receiver Operating Characteristic — Area Under the Curve)}

\newabbreviation{f1}{F1}{F1-Score: αρμονικός μέσος ακρίβειας και ανάκλησης}

\newabbreviation{cv}{CV}{Σταυροειδής Επαλήθευση (Cross-Validation)}

\newabbreviation{tp}{TP}{Αληθώς Θετικό (True Positive)}

\newabbreviation{fp}{FP}{Ψευδώς Θετικό (False Positive)}

%% — Signal Processing & Acoustic Features ————————————————————————————————————

\newabbreviation{mfcc}{MFCC}{Mel-Frequency Cepstral Coefficients}

\newabbreviation{hnr}{HNR}{Αρμονική προς Θορυβώδη Αναλογία
    (Harmonic-to-Noise Ratio)}

\newabbreviation{lpc}{LPC}{Γραμμική Προγνωστική Κωδικοποίηση
    (Linear Predictive Coding)}

\newabbreviation{tqwt}{TQWT}{Μετασχηματισμός Κυματιδίου Ρυθμιζόμενου Παράγοντα Q
    (Tunable Q-factor Wavelet Transform)}

%% — Datasets & Repositories ——————————————————————————————————————————————————

\newabbreviation{mdvrkcl}{MDVR-KCL}{Mobile Device Voice Recordings at King's College London}

\newabbreviation{pdsf}{PDSF}{Χαρακτηριστικά Ομιλίας Νόσου Πάρκινσον
    (PD Speech Features dataset)}

\newabbreviation{uci}{UCI}{University of California, Irvine Machine Learning Repository}

%% — Reporting Standards ——————————————————————————————————————————————————————

\newabbreviation{tripodai}{TRIPOD+AI}{Transparent Reporting of a multivariable prediction model for Individual Prognosis Or Diagnosis + AI}

\newabbreviation{probastai}{PROBAST+AI}{Prediction model Risk Of Bias Assessment Tool + AI}


%% =============================================================================
%% PART B — TECHNICAL TERMS (type = main)
%% =============================================================================

%% — Machine Learning ————————————————————————————————————————————————————————

\newglossaryentry{crossvalidation}{
    name        = {cross-validation},
    description = {Μέθοδος αξιολόγησης μοντέλου μηχανικής μάθησης κατά την
        οποία τα δεδομένα χωρίζονται επανειλημμένα σε σύνολα εκπαίδευσης και
        ελέγχου. Επιτρέπει την εκτίμηση της γενικευσιμότητας του μοντέλου χωρίς
        διαρροή δεδομένων. Στην παρούσα εργασία, η εφαρμογή γίνεται σε επίπεδο
        υποκειμένου για το Dataset A (ομαδοποιημένη στρωματοποιημένη 5-fold CV)
        και σε επίπεδο δείγματος για το Dataset B (τυπική στρωματοποιημένη
        5-fold CV)}
}

\newglossaryentry{overfitting}{
    name        = {overfitting},
    description = {Υπερπροσαρμογή: φαινόμενο κατά το οποίο ένα μοντέλο
        μαθαίνει υπερβολικά καλά τα δεδομένα εκπαίδευσης, με αποτέλεσμα
        κακή απόδοση σε νέα δεδομένα (απώλεια γενικευσιμότητας)}
}

%% — Signal Processing & Acoustics ————————————————————————————————————————————

\newglossaryentry{jitter}{
    name        = {jitter},
    description = {Διαταραχή περιόδου (Jitter): μέτρο κύκλο-προς-κύκλο
        αστάθειας της θεμελιώδους περιόδου δόνησης των φωνητικών χορδών.
        Υψηλές τιμές jitter παρατηρούνται στην υποκινητική δυσαρθρία
        που συνδέεται με τη νόσο Πάρκινσον}
}

\newglossaryentry{shimmer}{
    name        = {shimmer},
    description = {Διαταραχή πλάτους (Shimmer): μέτρο κύκλο-προς-κύκλο
        αστάθειας του πλάτους της δόνησης των φωνητικών χορδών. Μαζί
        με το jitter, αποτελεί ένδειξη νευρομυϊκής αστάθειας στη
        νόσο Πάρκινσον}
}

%% — Clinical & Physiological Terms ——————————————————————————————————————————

\newglossaryentry{dysarthria}{
    name        = {dysarthria},
    description = {Δυσαρθρία: νευρολογική διαταραχή της κινητικής παραγωγής
        ομιλίας που επηρεάζει τον συντονισμό των φωνητικών μυών.
        Ο υποτύπος που σχετίζεται με τη νόσο Πάρκινσον ονομάζεται
        υποκινητική δυσαρθρία (HKD)}
}

\newglossaryentry{hypophonia}{
    name        = {hypophonia},
    description = {Υποφωνία: μειωμένη ένταση ή φωνητική έκφραση της ομιλίας,
        συχνά παρατηρείται στη νόσο Πάρκινσον ως αποτέλεσμα εξασθένησης
        των αναπνευστικών και φωνητικών μυών}
}

\newglossaryentry{monopitch}{
    name        = {monopitch},
    description = {Μονοτονικότητα: μειωμένη μεταβλητότητα της 
        θεμελιώδους συχνότητας (F0), που οδηγεί σε μονοτονική ομιλία.
        Αποτελεί χαρακτηριστικό της υποκινητικής δυσαρθρίας στη νόσο
        Πάρκινσον και επηρεάζει την προσωδία και την εκφραστικότητα}
}

\newglossaryentry{bradykinesia}{
    name        = {bradykinesia},
    description = {Βραδυκινησία: κλινικό χαρακτηριστικό της νόσου Πάρκινσον που
        χαρακτηρίζεται από επιβράδυνση των εκούσιων κινήσεων, επηρεάζοντας
        την άρθρωση και την προσωδία της ομιλίας}
}

